\documentclass[10pt,a4paper]{amsart}
\usepackage[margin=19mm]{geometry}
\usepackage{amsmath,amssymb,mathtools}
\usepackage[scr=rsfs,cal=euler]{mathalfa}
\usepackage{xfrac}
\usepackage[unicode,hidelinks,bookmarks=false]{hyperref}
\newcommand{\ie}{i.e.\ }
\newcommand{\cf}{cf.\ }
\newcommand{\resp}{respectively\ }
\newcommand{\Subsection}{\subsection}
\providecommand{\nobreakdash}{\nobreak\hbox{-}}
\setlength{\emergencystretch}{2em}
\multlinegap0pt
\usepackage{amssymb}
\usepackage{float}
\usepackage{tikz}
\usepackage{dsfont}
\usepackage{xcolor}
\usepackage{bm}
\newtheorem{lemma}{Lemma}
\title{An Interface Green's Function Framework for Complete Discrete $W^{1,\infty}$ Analysis of Discontinuous Galerkin Methods}
\author{Haitao Leng, Weifeng Qiu}
\date{Source: arXiv:2608.11559v4 (CC BY 4.0)}
\begin{document}
\maketitle

\noindent\emph{Source and licence.} An Interface Green's Function Framework for Complete Discrete $W^{1,\infty}$ Analysis of Discontinuous Galerkin Methods. Version arXiv:2608.11559v4 (CC BY 4.0), distributed by arXiv under the Creative Commons Attribution 4.0 International licence, http://creativecommons.org/licenses/by/4.0/, as recorded in the arXiv licence metadata for this exact version and shown on the abstract page https://arxiv.org/abs/2608.11559v4. Full licence notice: \texttt{licenses/arxiv-2608.11559-CC-BY-4.0.txt}.\par\smallskip

\noindent\textit{Editorial scope.} Counted unit: the article's lemma rG-bi-lee (source0.tex lines 1786--1795). The article's complete original proof of it is delivered with it (source0.tex lines 1796--1851, one \texttt{proof} environment of 56 lines). No mathematics is written, repaired or re-derived: the statement and the proof are the article's own lines, in the article's own notation, and no retained line is substituted or re-flowed.  Retained uncounted as the source-local context the proof needs: lines 933--935; lines 1742--1744; lines 1749--1752; lines 1770--1772.  Nothing was omitted from any retained span, so the package carries no omission record.  No retained line contains a citation, so the wrapper carries no bibliography and no bibliography entry.  No retained line contains a figure environment, an image-inclusion command or drawing code, so no image is delivered and the package declares no asset dependency.  Editorial formatting only, in addition: an AMS \texttt{amsart} preamble that restates the article's own declarations and macros used by the retained lines, namely the environments \texttt{lemma} on separate counters, together with the standard packages those lines need.  The retained environments print in this document as Lemma 1 in order of appearance, whereas the article prints the same items with the numbering of its own full document.  The complete native source of the article as distributed in the arXiv e-print for this exact version is bundled unchanged as \texttt{sources/207428/source0.tex}.
\begin{align}
\chi(z)=\int_{K_z}\delta_{h,z}(x)\chi(x)dx\quad \forall \chi\in \mathcal{P}^s(K_z),\quad \int_{K_z}\delta_{h,z}(x)dx=1,\label{ddd:1}
\end{align}

\begin{align}
|D_z^{\alpha}D_x^{\beta}G(z,x)|\leq C|x-z|^{2-|\alpha|-|\beta|},\quad \forall 1\leq |\alpha|,|\beta|\leq 2,~3\leq |\alpha|+|\beta|\leq 4,\label{pe-G-bi}
\end{align}

\begin{align}
\Delta^2g_{{\rm if}}(x_{\ast},x)=\nabla\cdot({\bm e}_i\delta_{h,K^+_{\ast}})-\nabla\cdot({\bm e}_i\delta_{h,K^-_{\ast}})\quad {\rm in}~\Omega,\quad
g_{{\rm if}}=\nabla g_{{\rm if}}\cdot{\bm n}=0\quad {\rm on}~\partial\Omega,\label{rG-bi}
\end{align}

\begin{align}
\Delta^2g_{{\rm bf}}(x^{\ast},x)=\nabla\cdot({\bm e}_i\delta_{h,x^{\ast}})\quad {\rm in}~\Omega,\quad g_{{\rm bf}}=\nabla g_{{\rm bf}}\cdot{\bm n}=0\quad {\rm on}~\partial\Omega,\label{rG1-bi}
\end{align}

\begin{lemma}\label{rG-bi-lee}
Let $g_{{\rm if}}$ and $g_{{\rm bf}}$ be the solutions of problems (\ref{rG-bi}) and (\ref{rG1-bi}), then we have
\begin{align}
|D^2g_{{\rm if}}|_{L^{\infty}(\Omega_j^2(x_{\ast}))}+|D^2g_{{\rm bf}}|_{L^{\infty}(\Omega_j^2(x^{\ast}))}\leq Chd_j^{-2},
\end{align}
and
\begin{align}
|D^2g_{{\rm if}}|_{L^2(\Omega_j^2(x_{\ast}))}+|D^2g_{{\rm bf}}|_{L^2(\Omega_j^2(x^{\ast}))}\leq Chd_j^{-1}.
\end{align}
\end{lemma}

\begin{proof}
For any $x\in \Omega_j^2(x_{\ast})$, using the definitions of the Green's function $G$ and the Dirac delta function $\delta_x$ we have
\begin{align}
g_{{\rm if}}(x_{\ast},x)=&(\delta_x(y),g_{{\rm if}}(x_{\ast},y))=(D^2G(x,y),D^2g_{{\rm if}}(x_{\ast},y))\nonumber\\
=&(\nabla\cdot({\bm e}_i\delta_{h,K^+_{\ast}})-\nabla\cdot({\bm e}_i\delta_{h,K^-_{\ast}}),G(x,y))_{\mathcal{T}_h}\nonumber\\
=&({\bm e}_i\delta_{h,K^-_{\ast}},\nabla G(x,y))_{K^-_{\ast}}-({\bm e}_i\delta_{h,K^+_{\ast}},\nabla G(x,y))_{K^+_{\ast}}.\nonumber
\end{align}
Therefore,
\begin{align}
\partial_{x_ix_j} g_{{\rm if}}(x_{\ast},x)=&\int_{K^-_{\ast}}({\bm e}_i\delta_{h,K^-_{\ast}})(y)\cdot(\partial_{x_ix_j}\nabla_y G(x,y)) dy\label{rG-bi-lee-proof:1}\\
&-\int_{K^+_{\ast}}({\bm e}_i\delta_{h,K^+_{\ast}})(y)\cdot(\partial_{x_ix_j}\nabla_y G(x,y)) dy.\nonumber
\end{align}
Since $\partial_{x_ix_j}\nabla_y G(x,y)$, as a function of the second variable $y$, has the first order derivative in $B(x_{\ast},2h)\cap\Omega$ (see (\ref{pe-G-bi})), using Taylor's expansion
we know that there are two constants $t_3,t_4\in (0,1)$ depending on $y$ such that
\begin{align*}
\partial_{x_ix_j}\nabla_y G(x,y)=\partial_{x_ix_j}\nabla_y G(x,x_{\ast})+\partial_{x_ix_j}D^2_y G(x,y+t_3(x_{\ast}-y))(y-x_{\ast}),\quad \forall y\in K_{\ast}^-,\\
\partial_{x_ix_j}\nabla_y G(x,y)=\partial_{x_ix_j}\nabla_y G(x,x_{\ast})+\partial_{x_ix_j}D^2_y G(x,y+t_4(x_{\ast}-y))(y-x_{\ast}),\quad \forall y\in K_{\ast}^+,
\end{align*}
where $y+t_3(x_{\ast}-y)$ for $y\in K_{\ast}^+$ and $y+t_4(x_{\ast}-y)$ for $y\in K_{\ast}^-$ belong to $B(x_{\ast},2h)\cap\Omega$ since the domain $\Omega$ is convex. Moreover, according to (\ref{ddd:1}) we have
\begin{align*}
({\bm e}_i\delta_{h,K_{\ast}^+}(y),\partial_{x_ix_j}\nabla_y G(x,x_{\ast}))_{K_{\ast}^+}=({\bm e}_i\delta_{h,K_{\ast}^-}(y),\partial_{x_ix_j}\nabla_y G(x,x_{\ast}))_{K_{\ast}^-}.
\end{align*}
 Then, through (\ref{pe-G-bi}), (\ref{rG-bi-lee-proof:1}) and the above three identities, we obtain
\begin{align*}
\left|\partial_{x_ix_j} g_{{\rm if}}(x_{\ast},x)\right|\leq&\int_{K^-_{\ast}}\left|\delta_{h,K^-_{\ast}}\right|\left|\partial_{x_ix_j}D^2_y G(x,y+t_3(x_{\ast}-y))\right|\left|y-x_{\ast}\right| dy\\
&+\int_{K^+_{\ast}}\left|\delta_{h,K^+_{\ast}}\right|\left|\partial_{x_ix_j}D^2_y G(x,y+t_4(x_{\ast}-y))\right|\left|y-x_{\ast}\right| dy\\
\leq&C\int_{K^-_{\ast}}h\left|\delta_{h,K^-_{\ast}}\right|\left|t_3(x_{\ast}-x)+(1-t_3)(y-x)\right|^{-2} dy\\
&+C\int_{K^+_{\ast}}h\left|\delta_{h,K^+_{\ast}}\right|\left|t_4(x_{\ast}-x)+(1-t_4)(y-x)\right|^{-2} dy\leq Chd_j^{-2},
\end{align*}
and
\begin{align*}
|D^2 g_{{\rm if}}|_{L^2(\Omega_j^2(x_{\ast}))}\leq Chd_j^{-1}.
\end{align*}

For any $x\in\Omega_j^2(x^{\ast})$, similar to (\ref{rG-bi-lee-proof:1}), we have
\begin{align}
\partial_{x_ix_j}g_{{\rm bf}}(x^{\ast},x)=-\int_{K^{\ast}}({\bm e}_i\delta_{h,x^{\ast}})(y)\cdot(\partial_{x_ix_j}\nabla_yG(x,y)) dy.\label{rG-bi-lee-proof:2}
\end{align}
Since $\partial_{x_ix_j}\nabla_yG(x,y)$, as a function of the second variable $y$, has the first order derivative in $B(x^{\ast},2h)\cap\Omega$ (see (\ref{pe-G-bi})),
Taylor's expansion, together with $\partial_{x_ix_j}\nabla_yG(x,x^{\ast})=0$ for any $x\in\Omega_j^2(x^{\ast})$, implies that there exists a constant
$m_2\in(0,1)$ depending on $y$ such that
\begin{align*}
\partial_{x_ix_j}\nabla_yG(x,y)=\partial_{x_ix_j}D_y^2G(x,y+m_2(x^{\ast}-y))(y-x^{\ast}),\quad \forall y\in K^{\ast},
\end{align*}
where $y+m_2(x^{\ast}-y)\in B(x^{\ast},2h)\cap\Omega$ since the domain $\Omega$ is convex.
Thus, (\ref{pe-G-bi}), (\ref{rG-bi-lee-proof:2}) and the above identity yield
\begin{align*}
\left|\partial_{x_ix_j}g_{{\rm bf}}(x^{\ast},x)\right|\leq&\int_{K^{\ast}}\left|\delta_{h,x^{\ast}}\right|\left|\partial_{x_ix_j}D_y^2G(x,y+m_2(x^{\ast}-y))\right|\left|y-x^{\ast}\right|dy\\
\leq&C\int_{K^{\ast}}h\left|\delta_{h,x^{\ast}}\right|\left|m_2(x-x^{\ast})+(1-m_2)(x-y)\right|^{-2} dy\leq Chd_j^{-2},
\end{align*}
and
\begin{align*}
|D^2 g_{{\rm bf}}|_{L^2(\Omega_j^2(x^{\ast}))}\leq Chd_j^{-1}.
\end{align*}
We conclude the proof.
\end{proof}

\end{document}
